\documentclass{article}
\begin{document}

\section{Overview}

A collection of tools for professors and students. These are developed as modules that can be developed independently by internship students over multiple semester. These modules can then be accessed via the personal website.

\section{The project}

\subsection{Contributors section}

It would be nice to display a list of people who have worked on the website or modules along with their contribution. We shall allow the option for the contributors to opt-out.

Aprroach (1): use a simple Excel sheet to track and provide source of truth for the website. Down side is not very scalable. Approach (2): create a unified inferface. Not much funding for hosting. The solution is to build an offline application that uses GitHub API and manage the Excel sheets, which solves the scalability issue of the first approach.

\subsection{Demonstrators section}

Similar to contributors.

\subsection{Theses section}

Similar to contributors.

\subsection{Automated publication section}

Using MTMT API, retrieve the publication data from MTMT. We will to modift the website and researve a route as special route. This special route will display basic data but it will mainly fetch data and populate a shell, which are the section for the publication.

\subsection{Student case and request management}

Students would navigate a decision tree which would allow the student to submit their request for letter of reference. After submission, a professor can review then, accept, edit, or reject the request.

\subsection{Automated Teams question and template answer}


\subsection{LaTeX template bulk editor}

\subsection{LaTeX exam quesetion generator from question bank}

\subsection{Dedicated page for demonstrator information}

Shows how many slots will be open for demonstator next semester, responsibility and option for demostrator feedback.

<<<>>>

\end{document}